\documentclass[11pt,a4paper]{article}
\usepackage[T1]{fontenc}
\usepackage{lmodern}
\usepackage{amsmath,amssymb,amsthm,mathtools,mathrsfs}
\usepackage[margin=25mm]{geometry}
\usepackage{microtype,booktabs,longtable,array,enumitem,hyperref,xurl,graphicx}
\hypersetup{colorlinks=true,linkcolor=blue,citecolor=blue,urlcolor=blue,
 pdftitle={LM4: Finite group actions, loop structures, and profinite extensions},
 pdfauthor={AI-produced by GPT-6 Astra (OpenAI, ChatGPT Pro mode)}}
\setlength{\emergencystretch}{1.5em}
\setlength{\parskip}{3pt}
\newcommand{\Z}{\mathbb Z}\newcommand{\Q}{\mathbb Q}
\newcommand{\R}{\mathbb R}\newcommand{\C}{\mathbb C}\newcommand{\F}{\mathbb F}
\DeclareMathOperator{\Aut}{Aut}\DeclareMathOperator{\Out}{Out}
\DeclareMathOperator{\Inn}{Inn}\DeclareMathOperator{\Hom}{Hom}
\DeclareMathOperator{\id}{id}\DeclareMathOperator{\im}{im}
\DeclareMathOperator{\coker}{coker}\DeclareMathOperator{\Spec}{Spec}
\DeclareMathOperator{\Gal}{Gal}\DeclareMathOperator{\Map}{Map}
\DeclareMathOperator{\Homeo}{Homeo}\DeclareMathOperator{\Diff}{Diff}
\DeclareMathOperator{\rank}{rank}\DeclareMathOperator{\Sec}{Sec}
\DeclareMathOperator{\GL}{GL}\DeclareMathOperator{\PGL}{PGL}
\newenvironment{claim}[3]{\par\medskip\phantomsection\label{cl:#1}%
 \noindent\textbf{#1. [#2] #3.}\ \ignorespaces}{\par\medskip}
\newcommand{\passage}[1]{\par\noindent{\footnotesize\emph{Manuscript:} #1.}\par}
\begin{document}
\begin{center}
{\Large\bfseries Finite group actions, loop structures,\par
and profinite extensions\par}\smallskip
{\large A mathematical development of \emph{La Longue Marche}: LM4\par}\medskip
\end{center}
\noindent AI-produced by GPT-6 Astra (OpenAI, ChatGPT Pro mode). No priority claimed.
\par\noindent 13 November 2026

\noindent\textbf{Working manuscript: INCOMPLETE.}
The numbered statements distinguish proved assertions, complete counterexamples,
source theorems with their reductions, and the exact obligations still under
investigation. This working record makes no assertion of completion of the task.

\smallskip
\noindent\textbf{Notation and its manuscript concordance.}
Local generators, base points, coefficients and comparison maps are defined at
their use. The table fixes the global symbols and the meaning of each variant.

{\small\renewcommand{\arraystretch}{1.13}\noindent
\begin{tabular}{@{}p{.16\textwidth}p{.49\textwidth}p{.29\textwidth}@{}}
\toprule
Our notation & Meaning and Grothendieck's notation & Introducing passage / scans\\
\midrule
$X,S,U=X-S$ & Compact surface, marked set, punctured surface; $X,S,U$ & \S12: 91--92; \S19: 204\\
$(g,\nu),I$ & Genus, number and index set of punctures; $(g,\nu),I$ & \S13 bis: 115--116\\
$\pi,\pi_0,P$ & Discrete surface group, chosen discretization, its completion; $\pi,\pi_0,\pi$ & \S12: 91; \S28: 352\\
$L_i,\kappa_i$ & Indexed peripheral subgroup and exterior source map; $L_i,\kappa_i$ & \S12: 91--99\\
$T=T_\pi$ & Covariant orientation line; dual of the earlier cohomological line & 91; 116; 130\\
$A,A^0$ & Surface homeomorphism group and identity component; $A_{g,\nu},A^0$ & \S19: 204--205\\
$M,\Gamma_{g,\nu}$ & Geometric mapping class group; $\mathcal T_{g,\nu},\Gamma_{g,\nu}$ & 193--203; 204--205\\
$!,+$ & Pointwise marked-point preservation; orientation preservation & 204--205\\
$G,\Gamma,E$ & Finite action, operator group, lift extension; his $\Gamma,G,E$ as locally specified & \S15: 137--139; 307--312\\
$\Phi,\mu$ & Maximal partial-section data and torsion multiplier; $\Phi,\mu$ & 150--160; 197--203\\
$SE,S\Gamma$ & Special actual and exterior automorphism groups; $SE,S\Gamma$ & \S25: 296--300\\
$B,\overline M$ & $\Out_{\rm lac}(P)$ and closure of a specified geometric group; $\widehat{\widehat{\mathcal T}},\widehat{\mathcal T}$ where identified & 318--320; 352--358\\
$\operatorname{Discr},\operatorname{Prediscr}$ & Actual, outer, or closed-stabilizer cosets as explicitly specified & \S28: 352--358\\
$GT_0,\lambda,f$ & The normalized centralizer of $S_3$, multiplier and commutator parameter & Two-group problem: 312, 316--317\\
\bottomrule
\end{tabular}}

\clearpage
\section*{Reading the proof record}
\addcontentsline{toc}{section}{Reading the proof record}
We work in the order A--F of the supplied task. A group with loops retains its
indexed exterior peripheral maps and its orientation line; this includes the
exceptional cases where the image subgroups alone lose that data. All
homomorphisms between profinite or pro-$p$ groups are continuous. A section
retains its prescribed quotient labels, and its equivalence is conjugacy by
the named kernel. Every passage citation gives the manuscript section, folio,
and scan; a folio followed by ``verso'' names its reverse side. A transcription
alternative is retained when mathematical evidence bears on it. Our corrections
are assertions about the displayed mathematics, not declarations that a
doubtful reading has been settled.

The results register contains the full numbered assertions and dependencies.
The readings register records the editorial consequences. The report gives
the obligation-by-obligation state, source checks and computation provenance.
An ``Open'' label records an unfinished obligation in this working manuscript;
it does not replace a requested proof. A ``Known'' label includes the
derivation from the precise source theorem. The code uses exact finite
arithmetic, and the proofs identify what the computations certify and what
requires a separate infinite-group argument.

\tableofcontents
\clearpage
